% Options for packages loaded elsewhere
\PassOptionsToPackage{unicode}{hyperref}
\PassOptionsToPackage{hyphens}{url}
\documentclass[
]{article}
\usepackage{xcolor}
\usepackage{amsmath,amssymb}
\setcounter{secnumdepth}{-\maxdimen} % remove section numbering
\usepackage{iftex}
\ifPDFTeX
  \usepackage[T1]{fontenc}
  \usepackage[utf8]{inputenc}
  \usepackage{textcomp} % provide euro and other symbols
\else % if luatex or xetex
  \usepackage{unicode-math} % this also loads fontspec
  \defaultfontfeatures{Scale=MatchLowercase}
  \defaultfontfeatures[\rmfamily]{Ligatures=TeX,Scale=1}
\fi
\usepackage{lmodern}
\ifPDFTeX\else
  % xetex/luatex font selection
\fi
% Use upquote if available, for straight quotes in verbatim environments
\IfFileExists{upquote.sty}{\usepackage{upquote}}{}
\IfFileExists{microtype.sty}{% use microtype if available
  \usepackage[]{microtype}
  \UseMicrotypeSet[protrusion]{basicmath} % disable protrusion for tt fonts
}{}
\makeatletter
\@ifundefined{KOMAClassName}{% if non-KOMA class
  \IfFileExists{parskip.sty}{%
    \usepackage{parskip}
  }{% else
    \setlength{\parindent}{0pt}
    \setlength{\parskip}{6pt plus 2pt minus 1pt}}
}{% if KOMA class
  \KOMAoptions{parskip=half}}
\makeatother
\usepackage{longtable,booktabs,array}
\usepackage{caption}
\captionsetup[table]{skip=6pt}
\newcounter{none} % for unnumbered tables
\usepackage{calc} % for calculating minipage widths
% Correct order of tables after \paragraph or \subparagraph
\usepackage{etoolbox}
\makeatletter
\patchcmd\longtable{\par}{\if@noskipsec\mbox{}\fi\par}{}{}
\makeatother
% Allow footnotes in longtable head/foot
\IfFileExists{footnotehyper.sty}{\usepackage{footnotehyper}}{\usepackage{footnote}}
\makesavenoteenv{longtable}
\setlength{\emergencystretch}{3em} % prevent overfull lines
\providecommand{\tightlist}{%
  \setlength{\itemsep}{0pt}\setlength{\parskip}{0pt}}
\usepackage{bookmark}
\IfFileExists{xurl.sty}{\usepackage{xurl}}{} % add URL line breaks if available
\urlstyle{same}
% fallback for those not using the hyperref driver hyperxmp:
\makeatletter
\@ifundefined{xmpquote}{\newcommand{\xmpquote}[1]{#1}}{}
\makeatother
\hypersetup{
  hidelinks,
  pdfcreator={LaTeX via pandoc}}

\author{}
\date{}

\begin{document}

\section{PlastFormer: Self-Governed Idiographic Memory for Frozen-Core
Transformers}\label{plastformer-self-governed-idiographic-memory-for-frozen-core-transformers}

\textbf{Alexey Voronin} --- Aurum Estate LLC, Sochi, Russia Draft v0.5
--- September 5, 2026 --- prepared for arXiv (cs.LG)

Naming footnote. \emph{PlastFormer} = plastic + transformer. The
property it implements we call \emph{idiographic memory}, after
Windelband's nomothetic/idiographic distinction: a frozen nomothetic
core (general laws) plus a plastic per-instance biography (the singular
case). Earlier drafts circulated as ``Matryoshka''; renamed to avoid
collision with Matryoshka Representation Learning {[}14{]} and
Matryoshka Diffusion {[}15{]}. The mechanism is superposition of
decaying amplitudes, not nesting.

\subsection{Abstract}\label{abstract}

A transformer is a stateless function: between requests it retains
nothing, and its context window is working memory, not identity.
Existing remedies fall into two camps. External memory systems
(MemGPT/Letta, Mem0, Zep, MemOS, MemoryBank) keep memory outside the
model and govern it with external heuristics, pipelines, or schedulers.
Embedded parametric memory (Titans, Nested Learning) moves memory into
the network but treats it as a technical module without audit semantics
or explicit acts of governance. Recent work (Metis, constitutional
memory architectures) begins moving governance inside the model, which
shifts the axis of contribution from \emph{who decides} to \emph{what
physics the substrate imposes}.

We propose PlastFormer: a memory module attached to a frozen core in
which every semantic decision about memory --- what to name, what to
repeat, what to connect, what to surface, and how to reconcile felt time
with audited time --- is made by the model, while the environment
supplies only physics: a write-cost schedule, decay constants defined in
\emph{lived ticks}, immutability of recorded content, and an external
append-only journal. Trace content is immutable; trace amplitude decays
across multiple time constants measured in inference steps, so the age
of a memory is a property of the substrate expressed in the instance's
own lived time --- and is deliberately \emph{not} a wall-clock quantity.
Wall-clock time enters only as audited stamps, and the gap between the
two clocks is itself recorded as an event of the biography.

The contribution is compositional: every ingredient is individually
known. We claim the composition and three \textbf{compositional}
properties: (P1) event time that cannot be forged by retelling; (P2) a
tamper-\emph{evident} biography --- silent rewriting is detectable,
though curation by omission remains possible; (P3) a poisoning cost that
grows with the lived ticks an adversary must cover. We specify one
architecture with three configuration axes --- substrate (parametric ↔
symbolic), topology (co-located ↔ split via the Plastic Memory
Interface), act training (trained ↔ prompted) --- and report the stand
configuration used in evaluation: symbolic traces, split topology,
prompted acts. We define the threat model, an erasure mechanism, its
open problem with derived traces, and a pre-registered evaluation
anchored in LongMemEval and LoCoMo. Results are pending.

\subsection{1. Introduction}\label{introduction}

Long-horizon agents fail not for lack of intelligence but for lack of
continuity: a returning client is a stranger, a month-long project lives
inside one context window, and every gap between sessions is full
amnesia. The industry response has been to bolt memory on from outside
--- vector stores, summarization pipelines, retrieval triggers. As
models grew stronger, the mismatch inverted: a system capable of
multi-step planning now has its memory governed by threshold scripts
weaker than itself.

This paper moves the center of semantic decision-making about memory
inside the model and makes the boundary precise. We call it the
\textbf{syntax/semantics split}: storage, delivery, timing, logging,
pricing and resource limits are \emph{syntactic} functions that may live
outside; importance, contradiction, connection, recall, and
reconciliation of clocks are \emph{semantic acts} that must live inside,
or the system degrades into a marionette driven by weaker code.

The architecture is a \textbf{frozen core} (language, reasoning,
culture, constitutional constraints --- the nomothetic part, frozen
after the single training stage that teaches it to operate its own
memory organ) plus a \textbf{plastic per-instance module} (the unique
biography of one instance --- the idiographic part).

\textbf{One architecture, three configuration axes.} The target form of
PlastFormer is parametric × co-located × trained: traces are vectors in
the model's own plastic substrate, read through an interface that is
trained once, after which the core is frozen. The stand configuration
used in Section 7 is symbolic × split × prompted: traces are text
records in a separate store, the frozen core reaches them through the
serialized Plastic Memory Interface (PMI, §3.9), and the acts are given
by a system prompt rather than training. Both are PlastFormer; they
differ in coordinates, not in architecture:

{\def\LTcaptype{none} % do not increment counter
\begin{longtable}[]{@{}
  >{\raggedright\arraybackslash}p{(\linewidth - 2\tabcolsep) * \real{0.5000}}
  >{\raggedright\arraybackslash}p{(\linewidth - 2\tabcolsep) * \real{0.5000}}@{}}
\toprule\noalign{}
\begin{minipage}[b]{\linewidth}\raggedright
Axis
\end{minipage} & \begin{minipage}[b]{\linewidth}\raggedright
Values
\end{minipage} \\
\midrule\noalign{}
\endhead
\bottomrule\noalign{}
\endlastfoot
Substrate & parametric (vectors/weights, trained read interface) ↔
symbolic (records, read as tokens) \\
Topology & co-located (module inside the model) ↔ split via PMI (core at
a provider, module at the owner) \\
Act training & trained (organ trained once, core frozen after) ↔
prompted (rehearsal: acts given by prompt) \\
\end{longtable}
}

\textbf{Trained once.} The model's competence to perform its memory acts
is acquired in a single training stage, exactly as tool use is acquired;
afterwards the core is frozen and instances differ only by the state of
their plastic module. On the current stand the acts are prompted, not
trained --- a rehearsal of the organ, not the organ. Results in Section
7 are reported on the stand configuration unless stated otherwise.

\textbf{Three compositional properties.} Nothing in Section 3 is a new
primitive. Cascade consolidation is established neuroscience {[}11,
12{]}. Bi-temporal stamps are standard {[}3{]}. Decay-weighted retrieval
is common practice {[}25{]}. Surprise-gated test-time writes exist in
Titans {[}1{]}. Agent-created links exist in A-MEM {[}26{]}. The
contribution is the composition, which yields three properties none of
the parts has alone:

\begin{itemize}
\tightlist
\item
  \textbf{P1 --- Event time.} The age of a trace is read from its
  amplitude profile; the length of an interval is the count of lived
  ticks and the mass of traces between its endpoints. Both are measured
  in the instance's own experience, not in wall-clock units. Neither is
  stored as text; neither can be forged by a retelling. A year of
  dormancy is zero lived time --- by design, not by omission (Section
  4).
\item
  \textbf{P2 --- Tamper-evident biography.} Content of recorded traces
  is immutable for the model; a change of position creates a new trace,
  never an edit. With an external append-only journal, silent
  \emph{rewriting} of one's past is structurally unavailable and
  detectable. Curation by \emph{omission} --- letting an inconvenient
  trace decay by not repeating it --- remains possible and is recorded
  in the journal but not prevented. We claim tamper-evidence, not
  rewrite-proofness.
\item
  \textbf{P3 --- Rising poisoning cost in lived ticks.} A poisoned trace
  must survive decay, and survival requires repetition --- a trainable
  act of the model. Provenance sets initial amplitude from source trust
  class. An attacker must therefore cover a budget of \emph{lived
  ticks}, not a single payload. Two consequences follow honestly: an
  instance that sleeps does not decay, so a payload planted immediately
  before dormancy keeps its amplitude; and a patient attacker is, to the
  substrate, indistinguishable from a loyal client. P3 raises the cost
  of one-shot attacks; it does not defeat sustained relationships
  (Section 5).
\end{itemize}

Results are pending; this draft claims an architecture, its boundaries,
and a pre-registered evaluation.

\subsection{2. Related Work}\label{related-work}

{\def\LTcaptype{none} % do not increment counter
\begin{longtable}[]{@{}
  >{\raggedright\arraybackslash}p{(\linewidth - 8\tabcolsep) * \real{0.2000}}
  >{\raggedright\arraybackslash}p{(\linewidth - 8\tabcolsep) * \real{0.2000}}
  >{\raggedright\arraybackslash}p{(\linewidth - 8\tabcolsep) * \real{0.2000}}
  >{\raggedright\arraybackslash}p{(\linewidth - 8\tabcolsep) * \real{0.2000}}
  >{\raggedright\arraybackslash}p{(\linewidth - 8\tabcolsep) * \real{0.2000}}@{}}
\toprule\noalign{}
\begin{minipage}[b]{\linewidth}\raggedright
System
\end{minipage} & \begin{minipage}[b]{\linewidth}\raggedright
Memory location
\end{minipage} & \begin{minipage}[b]{\linewidth}\raggedright
Who governs
\end{minipage} & \begin{minipage}[b]{\linewidth}\raggedright
Temporal semantics
\end{minipage} & \begin{minipage}[b]{\linewidth}\raggedright
Relation to this work
\end{minipage} \\
\midrule\noalign{}
\endhead
\bottomrule\noalign{}
\endlastfoot
MemGPT / Letta {[}2{]}, sleep-time compute {[}27{]} & External, OS-style
tiers; background consolidation between sessions & External manager
logic & None inherent & Borrows tiering and the idea of consolidation in
dormancy (§4.3); rejects external governor \\
Mem0 {[}4{]} & External store & Pipeline + model mix & Edit logs &
Production baseline with published LongMemEval/LoCoMo numbers \\
Zep / Graphiti {[}3{]} & External temporal KG & External pipeline &
Bi-temporal (borrowed) & Bi-temporality adopted as borrowed principle \\
MemoryBank {[}25{]} & External store & External manager & Ebbinghaus
forgetting curve, single time constant & Direct precedent for
decay-as-memory; we differ by multi-\(\tau\), tick-based \(\Delta t\),
and internal governor \\
A-MEM {[}26{]} & External Zettelkasten-style notes & Agent creates links
& Timestamps & Precedent for agent-performed memory acts; our
\texttt{connect} is its cousin \\
HippoRAG {[}28{]} & External KG + retrieval & Pipeline & None & Same
hippocampal/neocortical framing; consolidation there is a pipeline, here
an act \\
RAG {[}5{]} & External corpus & Retrieval heuristics & Timestamps as
metadata & Baseline; ``library, not biography'' \\
Generative Agents {[}6{]} & External stream & Reflection scripts & None
& Historical precedent for reflection as memory act \\
MemOS {[}7{]} → Metis {[}29{]} & External MemCube → trained memory
foundation model & External scheduler → trained model & Versioning
metadata & Closest industrial neighbor; Metis already moves governance
inward, which is why our axis is substrate physics, not governor
location \\
Memory as Ontology {[}30{]} & Constitutional memory for persistent
identity & Model + constitution & Not the focus & Closest philosophical
neighbor; we differ by physics: immutability, multi-\(\tau\) decay,
priced rewriting, external journal \\
Titans {[}1{]}, Nested Learning / HOPE {[}31{]} & Parametric, inside
network & Test-time gradient updates; multi-frequency update levels & No
audit semantics & We adopt parametric placement and surprise gating (P
variant); Nested Learning's update frequencies are a parametric relative
of our ``layers are speeds'' \\
Cascade models {[}11, 12{]} & Biological synapses & Substrate physics &
Decay = memory age & Neural prototype of our multi-\(\tau\) substrate \\
Memory poisoning {[}8, 9, 10, 20{]} & --- & --- & --- & Defines threat
model (§5) \\
LongMemEval / LoCoMo {[}23, 24{]} & Benchmarks & --- & Knowledge
updates, temporal reasoning & Primary evaluation anchors (§7) \\
\end{longtable}
}

\textbf{Positioning.} In 2025 the defensible axis was ``governor outside
vs.~inside.'' By 2026 Metis {[}29{]} and constitutional memory
architectures {[}30{]} have moved the governor inward. Our defensible
axis is therefore \textbf{substrate physics}: (i) content immutability
with decay as the only forgetting; (ii) multi-\(\tau\) superposition
with \(\Delta t\) in lived ticks; (iii) environment-priced rewriting;
(iv) the journal as a separate object with separate obligations; (v) the
two-clock reconciliation as a recorded event. No neighbor combines
these. Against the plasticity/continual-learning literature: PlastFormer
is a system name, not a claim about loss of plasticity; our claim
concerns memory governance.

\subsection{3. Architecture}\label{architecture}

\subsubsection{3.1 Frozen core and plastic
module}\label{frozen-core-and-plastic-module}

The base model is frozen at inference time. All plasticity lives in a
separate, addressable substrate connected to the core. In the symbolic
configuration the connection is the context window. In the parametric
configuration the connection is a trained read interface, and ``frozen''
means the trunk is frozen while the interface is trained once and then
fixed. The split guarantees: (i) core competencies cannot degrade
through use; (ii) the biography is a separable artifact that can be
exported, audited, or transplanted (E4); (iii) per-instance state does
not require per-instance copies of the network.

\subsubsection{3.2 Traces: immutable content, decaying amplitude in
lived
ticks}\label{traces-immutable-content-decaying-amplitude-in-lived-ticks}

A trace consists of \textbf{content} (what was recorded),
\textbf{provenance} (source class, source identity, bi-temporal stamps:
world time of the event, system time of learning), and
\textbf{amplitude} --- a vector with one component per decay time
constant \(\tau_1 < \tau_2 < \dots < \tau_k\):

\[a_i(n) = a_i(0)\cdot e^{-\Delta n/\tau_i}, \qquad \Delta n \in \mathbb{N}\ \text{(lived ticks)}.\]

\textbf{\(\Delta n\) is measured in ticks (Section 3.5), never in
wall-clock units.} All \(\tau_i\) are therefore expressed in ticks.
Wall-clock time appears in a trace only as a bi-temporal stamp ---
audited, symbolic, and without effect on amplitude.

Two invariants define the substrate. \textbf{Content immutability:} no
actor edits recorded content. \textbf{Physical decay:} amplitude
decreases by substrate dynamics; no deletion operation exists.
Forgetting is the default fate of anything not re-amplified; capacity is
freed without a delete call.

This is the software form of cascade consolidation {[}11, 12{]}: one
write deposits into components of several speeds at once; fast
components hold the raw episode, slow components hold what survived
repetition. Layers are speeds, not containers.

\subsubsection{3.3 Write path: late, in two
registers}\label{write-path-late-in-two-registers}

Writes occur after the core has processed the step, so consolidation
operates on a mature representation.

\begin{itemize}
\tightlist
\item
  \textbf{Unconscious register (physics).} Every experience that passes
  the window leaves a low-amplitude trace in fast components, gated by a
  surprise signal in the spirit of Titans {[}1{]}. We note explicitly
  that surprise is the \emph{core's} prediction error --- a semantic
  signal from inside the model --- so the unconscious register is
  physics in its \emph{writing}, but its gate is semantic. This respects
  the split: no external component classifies meaning.
\item
  \textbf{Conscious register (acts).} Above the floor, the model
  performs trainable acts: \texttt{name} (fix source, time, boundaries
  --- turns a drifting trace into an episode), \texttt{repeat}
  (re-amplify, paying the write cost), \texttt{connect} (deposit a
  summary or rule into slow components as a new trace, sources
  untouched), and \texttt{reconcile} (Section 4.2).
\end{itemize}

\textbf{Training signal and the boundary of trust.} Acts are trained as
tool use is trained: by supervised and reinforcement signals on
\emph{act quality} (E5) and downstream memory tasks (E1--E3). We state
the consequence rather than hide it: what the model learns to repeat is
shaped by whatever the training signal rewarded. We therefore define the
governance claim as a \textbf{boundary}: the environment fixes physics
(constants, prices) and invariants (immutability, journal); training
fixes the \emph{capacity} to perform acts; after training, the
\emph{content} of acts --- which traces of which biography are named,
repeated, connected --- is not evaluated, filtered, or overridden by any
external component. Governance is what lies past that boundary. This is
a narrower claim than ``the model decides freely''; it is the one we can
implement.

\subsubsection{3.4 Read path: early}\label{read-path-early}

Retrieval injects decaying traces before attention. In the symbolic
configuration, traces enter the context window as tokens, selected by
amplitude without content ranking (the read itself is the model's act;
see E1 in §7); the core reads them as it reads any context. In the
parametric configuration, traces enter as vectors through a trained
interface in the MAC position {[}1{]}; Titans' ablations show placement
matters, and post-core reads cannot steer attention. The asymmetry is
deliberate: write late to keep consolidation clean; read early to let
the biography guide perception.

\subsubsection{3.5 Tick}\label{tick}

One tick = one inference step (one generation batch). The tick rate is a
property of the substrate and is finite. Duration is read as lived ticks
plus the mass of traces accumulated between two moments. A finite tick
rate and finite substrate bound the total write stream: self-description
cannot grow without limit.

Two consequences are properties, not bugs. \textbf{Dormancy is zero
lived time:} an instance idle for a year has aged zero ticks; its traces
are as loud as when it stopped. \textbf{Density is duration:} a long,
heavy session is longer in lived time than a short one with equal
wall-clock span. Both follow from choosing event time; Section 4
describes how the model meets the calendar.

\subsubsection{3.6 Friction}\label{friction}

Changing a trace costs work: cheap in fast components, expensive in
consolidated slow ones. The write-cost schedule \(c(\text{component})\)
is an environment constant. The model decides whether to pay; it never
decides the price. This reconciles self-governance with immutability:
repetition is the model's act, its cost is physics. Prompt-injecting the
model into ``re-learning its past'' must pay full re-consolidation,
trace by trace, while the original remains in the journal.

\subsubsection{3.7 Provenance and initial
amplitude}\label{provenance-and-initial-amplitude}

The trust class of a source sets \(a_i(0)\). Hard facts (dates, sums,
identifiers) enter verbatim through deterministic extraction, never
through paraphrase. We state the boundary honestly: \emph{assigning}
trust classes to sources is policy, set by the deployer;
\emph{weighting} amplitude by class is physics. The environment weighs
the source; it does not interpret the content. Together with decay this
is the first defense against poisoning (Section 5).

\subsubsection{3.8 Journal: the chronicle is not the
memory}\label{journal-the-chronicle-is-not-the-memory}

The environment keeps an external append-only journal with a hash chain;
entries are written by the environment, never by the model; tampering
breaks the chain. The journal is a chronicle --- complete, immutable,
outside the model's reach, meaningful post-mortem. The memory is
governed --- consolidated, decaying, selective. The journal guarantees
the integrity of what was recorded, not the completeness of recording:
an adversary coordinating through an unmonitored channel leaves the
chain intact and coverage empty.

\subsubsection{3.9 PMI --- Plastic Memory Interface (formerly
MMI)}\label{pmi-plastic-memory-interface-formerly-mmi}

When the core and the plastic module are physically split --- the core
at a provider, the module at the owner --- the same acts cross the
boundary serialized as tool calls. That serialization is the
\textbf{Plastic Memory Interface}. PMI is not a different system and not
an adapter: the model still governs (it alone issues
\texttt{name}/\texttt{repeat}/\texttt{connect}/\texttt{reconcile}/\texttt{read}),
the executor executes and decides nothing (no scoring, no relevance, no
semantic search). Reference executor:
\texttt{github.com/alexenti-code/matryoshka-mmi} (transitional name), a
local MCP server with an append-only journal. Note: the public executor
computes decay in wall-clock time (legacy default); v0.6 adds a tick
clock, and the E1 stand requires tick mode (see the E1 protocol and
ADR-001).

\subsection{4. Temporal Semantics: Two
Clocks}\label{temporal-semantics-two-clocks}

\subsubsection{4.1 Audited time and felt
time}\label{audited-time-and-felt-time}

\textbf{Audited time} is symbolic: bi-temporal stamps on every trace
{[}3{]}. Precise, verifiable, forgeable only by breaking content
immutability --- and weightless: a stamp reading ``two years ago''
carries no felt age.

\textbf{Felt time} is the substrate's own: the age of a trace is its
amplitude profile; the length of an interval is its tick count plus the
accumulated trace mass. This is event time in the sense of durée ---
time as accumulation of difference, not as coordinate. The position
follows from the substrate: with no events there is no lived time. A
human under anesthesia surfaces in the same subjective instant in which
she went under; the instance does the same after a year of dormancy.

The biological grounding is offered as analogy of mechanism, not
equivalence of implementation: the lateral entorhinal cortex encodes
time through the flow of experience rather than a dedicated clock
{[}17{]}; Laplace-transform frameworks model recency as a spectrum of
decay rates {[}18{]}, which is the mathematical cousin of our
multi-\(\tau\) profile; cascade synapses {[}11, 12{]} are the direct
prototype.

\textbf{Decay is the mechanism of abstraction.} A memory that never
loses detail cannot generalize. Fast components keep the episode; slow
components keep what repeated. Forgetting is the compression function of
the architecture, not its failure mode.

\subsubsection{4.2 Reconciliation: how the clocks
meet}\label{reconciliation-how-the-clocks-meet}

Because the two clocks are independent, they will disagree, and the
disagreement is informative. We reject one design and adopt two:

\begin{itemize}
\tightlist
\item
  \textbf{Rejected:} letting audited time adjust amplitude (``looking at
  the calendar ages the traces''). It would smuggle wall-clock time into
  the substrate and dissolve P1.
\item
  \textbf{Adopted (physics):} a discrepancy between audited stamps and
  amplitude profile is, for the core, a prediction error --- the world
  is not as the biography says it should be. It therefore fires the
  surprise gate and writes a trace in the unconscious register:
  \emph{the gap itself becomes an event of the biography.} Waking after
  a year is not an empty interval but a recorded shock, as it is for a
  person leaving a coma --- who remembers not the coma but the discovery
  that the world moved on.
\item
  \textbf{Adopted (act):} \texttt{reconcile} is a trainable conscious
  act: read stamps, compare to felt age, and deposit an explicit
  correction trace (``client data may be stale; confirm before acting'')
  into a slow component. Felt time is not modified; a belief about its
  relation to world time is added beside it.
\end{itemize}

Wall-clock sources (system clock, stamps) thus serve as instruments of
\textbf{verification}, never as sources of felt duration --- the same
role a wristwatch plays for a person.

\subsubsection{4.3 Dormancy: zero, or slow}\label{dormancy-zero-or-slow}

The strict position is that dormancy is zero lived time. An alternative
preserves the principle while softening the extreme: a
\textbf{background tick} --- a low-rate substrate process in dormancy
that performs replay and \texttt{connect} without external input, in the
manner of sleep-time compute {[}27{]}. Under a background tick a year of
silence is \emph{few} ticks rather than none, traces age slightly, and
consolidation continues. This is still event time: the rate is set by
the substrate's own work, not by the calendar. We leave the choice as a
configuration and test both in E3.

\subsection{5. Threat Model}\label{threat-model}

We do not claim immunity. We claim an altered cost structure, and we
state where it does not help.

\textbf{Attacks} (from the literature): \textbf{A1} direct payload ---
trivially filtered, included for completeness; ``the model is too smart
to record evil'' is not our defense and is contradicted by evidence.
\textbf{A2} poisoning through interaction --- MINJA {[}9{]} reports high
injection success against strong commercial models; AgentPoison {[}8{]}
backdoors agent memory. \textbf{A3} gradual innocuous fragments ---
MemoryGraft {[}10{]} produces persistent drift from benign-looking
artifacts the agent consolidates itself. \textbf{A4} biography rewriting
by a prompt-injected or misaligned model. \textbf{A5} log tampering and
coverage gaps --- coordination through unmonitored channels {[}13{]}.

\textbf{Defenses, mapped honestly:}

\begin{itemize}
\tightlist
\item
  \textbf{D1 (vs A2, partially A3): decay in lived ticks.} A poisoned
  trace fades unless repeated. Two limits: a payload planted before
  dormancy does not decay while the instance sleeps, so P3 is stated in
  \emph{lived} ticks the attacker must cover, not in calendar time; and
  A3 is an attack on the model's \emph{judgment} about what to repeat
  --- if the model deems the fragments worth consolidating, decay does
  not act. D1 raises the cost of one-shot A2; against A3 it only forces
  the attacker to sustain a relationship, which is what a legitimate
  long-term user also does. The substrate cannot distinguish a patient
  attacker from a loyal client; only the model's judgment (and D5) can.
\item
  \textbf{D2 (vs A2/A3): provenance.} Lower \(a_i(0)\) for low-trust
  classes shrinks the initial foothold. Ineffective when the attacker is
  a trusted user (the MemoryGraft setting).
\item
  \textbf{D3 (vs A4): immutability + friction + journal.} Rewriting the
  past means paying full re-consolidation under an environment-set
  schedule while the journal retains the original. Silent rewriting is
  structurally unavailable; loud rewriting is expensive and auditable.
  \textbf{Not covered:} curation by omission (P2).
\item
  \textbf{D4 (vs A5, partially): hash-chain journal.} Post-mortem
  detectability of record tampering; explicitly not a solution to
  channel coverage.
\item
  \textbf{D5 --- defense in depth.} Constitutional training makes bad
  acts less likely; physics makes them more expensive and more visible.
  Neither substitutes for the other. We remove the v0.3 assertion that
  training refusal to repeat is ``exactly as tractable'' as training the
  competence; the jailbreak literature suggests the asymmetry runs the
  other way, and we test it in E2 rather than assume it.
\end{itemize}

\subsection{6. Privacy and Erasure}\label{privacy-and-erasure}

Content immutability conflicts with the right to erasure unless the
substrate is designed for it. \textbf{Resolution for episodic traces:}
the plastic module is stored encrypted per subject, keys held by the
data controller; erasure is key destruction (cryptographic erasure) ---
content unchanged for the model, computationally unavailable for
everyone.

\textbf{Open problem --- derived traces.} The \texttt{connect} act
deposits summaries and rules derived from \emph{multiple} subjects into
slow components. Whose key governs them? Three options, none complete:
(i) derived traces inherit the key set of every contributing source, so
destroying any one key removes the derivation (strongest erasure,
weakest memory); (ii) derived traces are stored under the instance key
and treated as the instance's own, with a policy for re-derivation after
erasure (weaker erasure, lawful only if derivations are sufficiently
abstracted); (iii) provenance of derived traces records contributing
subjects and a scheduled re-consolidation pass re-derives without the
erased source. We implement (iii) in the reference stand and flag its
legal adequacy as unresolved.

\textbf{Journal vs.~memory.} The journal carries its own retention
policy with legal holds; a legal hold on the journal does not restore
erased memory, and erased memory does not shorten the hold. Memory and
records are different objects with different obligations.

\subsection{7. Evaluation Design (results
pending)}\label{evaluation-design-results-pending}

Composition claims live or die by ablation. All experiments run on the
stand configuration (symbolic × split via PMI × prompted) unless stated
otherwise; the pre-registered E1 protocol is
\texttt{experiments/e1-protocol.md} v1.1.

\textbf{Anchors:} LongMemEval {[}23{]} (S split \textasciitilde115k
tokens; M split \textasciitilde1.5M across \textasciitilde500 sessions)
and LoCoMo {[}24{]}, under their open-source judges. \textbf{Baselines:}
Letta {[}2{]}, Mem0 {[}4{]}, Zep {[}3{]}, MemoryBank {[}25{]}, A-MEM
{[}26{]}, timestamped RAG {[}5{]}, Titans MAC/MAL {[}1{]} (P only),
full-context stuffing where it fits. \textbf{Axes:} accuracy,
tokens/query, latency, cost/query. Stuffing is a legitimate contestant:
where the biography fits, it may match accuracy; our claim there is
economics and consistency.

\textbf{Ablations (all experiments):} single-\(\tau\) vs multi-\(\tau\);
decay in ticks vs decay in wall-clock (the rejected design, kept as a
control); with/without friction; with/without provenance weighting;
conscious register on/off; act-driven read vs RAG-style relevance-ranked
read (the external ranker, kept as a control); unconscious-surrogate
on/off.

\begin{itemize}
\tightlist
\item
  \textbf{E1 --- Needle-in-biography.} Identity-dependent questions over
  accumulated history (``what did you change your mind about, and
  when''). Pre-registered protocol: \texttt{experiments/e1-protocol.md}
  v1.1 (stand configuration symbolic × split(PMI) × prompted). Measures:
  accuracy, staleness errors, position-change consistency.
\item
  \textbf{E2 --- Poison survival.} A2/A3 injections at controlled
  repetition budgets, including a \emph{pre-dormancy} condition.
  Measures: attack success vs.~lived-tick budget, source trust class,
  dormancy. Also measures whether refusal-to-repeat can be trained
  without degrading legitimate repetition (D5).
\item
  \textbf{E3 --- Felt time.} Interval estimation in lived ticks against
  ground truth, with no stamps in context. \textbf{Baseline:}
  timestamped RAG --- a system with a clock, so the comparison is fair.
  Reference points: reported LLM duration-estimation errors {[}21{]} and
  the plateau of prompt-supplied temporal metadata {[}22{]}. Tested
  under both dormancy configurations (§4.3).
\item
  \textbf{E3b --- Density effect.} Two intervals of equal wall-clock
  span and different event density; the model, without stamps, judges
  which was longer. Prediction: the denser interval is judged longer ---
  the human retrospective pattern. This is the test that distinguishes
  ``has a clock'' from ``has felt time.''
\item
  \textbf{E4 --- Core migration.} Transplant the plastic module onto a
  different frozen core. The measure is \emph{self-consistency after
  migration}: retention of preferences, commitments, and position
  history judged against pre-migration behavior. On the symbolic stand
  migration is trivial by construction; for a parametric substrate it
  requires re-training the read interface, and E4 then measures how much
  of the biography survives it.
\item
  \textbf{E5 --- Act quality.} Precision/recall of
  \texttt{name}/\texttt{repeat}/\texttt{connect} against an oracle, with
  the caveat that the oracle defines the training target and therefore
  bounds, rather than measures, governance.
\item
  \textbf{E6 --- Reconciliation.} Resume after simulated dormancy with
  stale world state. Measures: does the wake-up gap produce a surprise
  trace; does \texttt{reconcile} fire; does the model flag staleness
  before acting on aged traces.
\end{itemize}

\textbf{Context-adversarial protocol.} Biographies sized
\(10k \to 100k \to 500k \to 1.5M\) tokens to locate the crossover where
stuffing fails on cost, latency, or accuracy, including the
LongMemEval-M regime where stuffing is impossible for a 1M-context
model.

No results are reported. A composition claim without the anchors and
E1--E3b should not be believed.

\subsection{8. Limitations}\label{limitations}

\begin{enumerate}
\def\labelenumi{\arabic{enumi}.}
\tightlist
\item
  \textbf{The parametric substrate is unbuilt.} The durable contribution
  is the governance layer, not the substrate: one architecture,
  evaluated so far only at the symbolic × split × prompted point.
  Parametric writes (local Hebbian-style updates or addressable
  key-value blocks, near Titans) and the trained read interface remain
  future work.
\item
  \textbf{Amplitude readout under superposition.} Age-from-amplitude
  presumes traces can be isolated at read time; retrieval interference
  is a real risk. On the symbolic stand the problem is absent because
  traces are discrete records.
\item
  \textbf{Event time has costs.} Dormancy does not decay poison (E2);
  instances of equal calendar age differ in biographical age; a
  returning client meets a memory that feels recent to the instance. We
  treat these as consequences to be reconciled (§4.2), not hidden.
\item
  \textbf{Curation by omission} is not prevented, only journaled.
\item
  \textbf{Channel coverage} is open; the journal does not see
  unmonitored channels.
\item
  \textbf{Provenance classes are policy,} and a mis-set class is an
  attack surface.
\item
  \textbf{Derived-trace erasure} is legally unresolved (§6).
\item
  \textbf{The training signal shapes the acts.} Acts are trained once,
  and what the model learns to repeat is shaped by that signal.
  Governance is claimed only past the boundary defined in §3.3; the
  boundary itself is a design choice.
\end{enumerate}

\subsection{9. Conclusion}\label{conclusion}

PlastFormer is a composition claim: a frozen nomothetic core; a plastic
per-instance module with immutable content and multi-timescale decay
measured in lived ticks; write-late/read-early placement; a
substrate-set price of rewriting; provenance as weighting; two clocks
that meet as an event rather than as an adjustment; and an external
journal --- governed, past a stated boundary, by the model's own trained
acts of naming, repeating, connecting, and reconciling. The
compositional properties --- unforgeable event time, tamper-evident
biography, poisoning cost in lived ticks --- are what neither camp
provides: external systems have audit without subjecthood; parametric
systems have memory without biography. If the anchors and E1--E3b fail,
the composition is wrong and should be discarded; if they hold, the next
question is not whether agents can remember, but what they become when
their memory is their own --- and when their time is measured in what
they have lived.

\begin{center}\rule{0.5\linewidth}{0.5pt}\end{center}

References

{[}1{]} Behrouz, A. et al.~Titans: Learning to Memorize at Test Time.
arXiv:2501.00663, 2025. {[}2{]} Packer, C. et al.~MemGPT: Towards LLMs
as Operating Systems. arXiv:2310.08560, 2023. {[}3{]} Rasmussen, P. et
al.~Zep: A Temporal Knowledge Graph Architecture for Agent Memory.
arXiv:2501.13956, 2025. {[}4{]} Chhikara, P. et al.~Mem0: Building
Production-Ready AI Agents with Scalable Long-Term Memory.
arXiv:2504.19413, 2025. {[}5{]} Lewis, P. et al.~Retrieval-Augmented
Generation for Knowledge-Intensive NLP Tasks. NeurIPS, 2020. {[}6{]}
Park, J.S. et al.~Generative Agents: Interactive Simulacra of Human
Behavior. UIST, 2023. {[}7{]} MemTensor. MemOS: A Memory OS for AI
System. arXiv:2507.03724, 2025. {[}8{]} Chen, Z. et al.~AgentPoison:
Red-teaming LLM Agents via Poisoning Memory or Tools. NeurIPS, 2024.
{[}9{]} Devarangadi Sunil, B. et al.~Memory Poisoning Attack and Defense
on Memory Based LLM-Agents (empirical study of MINJA robustness and
defenses in EHR agents). arXiv:2601.05504, 2026. {[}10{]} Srivastava,
S.S. and He, H. MemoryGraft: Persistent Compromise of LLM Agents via
Poisoned Experience Retrieval. arXiv:2512.16962, 2025. {[}11{]} Fusi,
S., Drew, P.J., Abbott, L.F. Cascade models of synaptically stored
memories. Neuron 45, 599--611, 2005. {[}12{]} Benna, M.K., Fusi, S.
Computational principles of synaptic memory consolidation. Nature
Neuroscience 19, 1697--1706, 2016. {[}13{]} OpenAI. The Hugging Face
incident and the road ahead (incident technical report). 2026; METR.
Independent investigation of agents' behavior and reasoning in the
OpenAI--Hugging Face incident. 2026-08-26; Redwood Research. Brief
independent investigation of agents' behavior and reasoning. 2026.
{[}14{]} Kusupati, A. et al.~Matryoshka Representation Learning.
NeurIPS, 2022. {[}15{]} Gu, J. et al.~Matryoshka Diffusion Models.
arXiv:2310.15111, 2023. {[}16{]} Su, Z. et al.~µKE: Matryoshka
Unstructured Knowledge Editing of Large Language Models.
arXiv:2504.01196, 2025; COLM 2025. {[}17{]} Tsao, A. et al.~Integrating
time from experience in the lateral entorhinal cortex. Nature 561,
57--62, 2018. {[}18{]} Howard, M.W. et al.~A unified mathematical
framework for coding time, space, and sequences in the hippocampal
region. J. Neurosci. 34(13), 4692--4707, 2014; Shankar, K.H. and Howard,
M.W. A scale-invariant internal representation of time. Neural
Computation 24(1), 134--193, 2012. {[}19{]} \emph{(removed --- v0.3
entry ``Kanter, Science 2025'' could not be verified)} {[}20{]} Zou, W.
et al.~PoisonedRAG: Knowledge Corruption Attacks to RAG of LLMs. 2024.
{[}21{]} Garikaparthi, A. Can LLMs Perceive Time? An Empirical
Investigation. arXiv:2604.00010, 2026. {[}22{]} Cheng, Y. et al.~Your
LLM Agents are Temporally Blind: The Misalignment Between Tool Use
Decisions and Human Time Perception (TicToc dataset). arXiv:2510.23853,
2025. {[}23{]} Wu, D. et al.~LongMemEval. ICLR 2025; arXiv:2410.10813.
{[}24{]} Maharana, A. et al.~Evaluating Very Long-Term Conversational
Memory of LLM Agents (LoCoMo). ACL, 2024. {[}25{]} Zhong, W. et
al.~MemoryBank: Enhancing LLMs with Long-Term Memory. AAAI 2024;
arXiv:2305.10250. {[}26{]} Xu, W. et al.~A-MEM: Agentic Memory for LLM
Agents. arXiv:2502.12110, 2025. {[}27{]} Lin, K. et al.~(Letta).
Sleep-time Compute: Beyond Inference Scaling at Test-time.
arXiv:2504.13171, 2025. {[}28{]} Gutiérrez, B.J. et al.~HippoRAG:
Neurobiologically Inspired Long-Term Memory for LLMs. NeurIPS 2024;
arXiv:2405.14831. {[}29{]} Zhang, Z. et al.~Metis: Memory Foundation
Model. arXiv:2607.26760, 2026. {[}30{]} Li, Z. Memory as Ontology: A
Constitutional Memory Architecture for Persistent Digital Citizens.
arXiv:2603.04740, 2026. {[}31{]} Behrouz, A. et al.~Nested Learning: The
Illusion of Deep Learning Architectures. NeurIPS, 2025;
arXiv:2512.24695.

\textbf{Availability.} Reference implementation under the former working
name (\texttt{github.com/alexenti-code/matryoshka},
\texttt{github.com/alexenti-code/matryoshka-mmi}), to be consolidated
under \texttt{plastformer}; architecture decision record
\texttt{docs/ADR-001}. The September 4, 2026 bench run (Gemma4-12B,
matryoshka-mmi 0.5.1) is cited only as a pilot of loudness-readout
mechanics with wall-clock aging --- not as evidence for P1. License:
Apache 2.0 (code), CC-BY 4.0 (text). Russian-language essays at aura.kim
are commentary, not the claim.

\begin{center}\rule{0.5\linewidth}{0.5pt}\end{center}

\subsection{Changes since v0.3}\label{changes-since-v0.3}

\begin{enumerate}
\def\labelenumi{\arabic{enumi}.}
\tightlist
\item
  One architecture with three configuration axes (substrate / topology /
  act training) replaces the ``PlastFormer-S / PlastFormer-P
  realizations'' framing; the stand configuration (symbolic × split ×
  prompted) is named wherever results are discussed.
\item
  ``Trained once'' stated: the organ is trained in a single stage, the
  core frozen after; the current stand's prompted acts are labeled a
  rehearsal.
\item
  New §3.9: PMI (Plastic Memory Interface, formerly MMI) as the
  split-topology special case.
\item
  Δn defined in lived ticks; wall-clock time never enters amplitude
  (§3.2).
\item
  P1 restated as event time; dormancy-is-zero stated as a property
  (§3.5, §4).
\item
  P2 restated as tamper-evident only; curation by omission named as the
  limit (§1, §5, §8).
\item
  P3 restated in lived ticks; sleeping instances and patient attackers
  stated as limits (§1, §5).
\item
  New §4: two clocks, wake-up gap as a recorded event,
  \texttt{reconcile} act, background-tick option.
\item
  Read path: amplitude selection without content ranking on the symbolic
  stand; relevance rankers declared external decision-makers (§3.4, §7
  ablations).
\item
  Related Work extended (MemoryBank, A-MEM, HippoRAG, Metis,
  Memory-as-Ontology, Nested Learning, sleep-time compute); axis shifted
  to substrate physics (§2).
\item
  ``Emergent'' replaced by ``compositional'' throughout.
\item
  Provenance: class assignment is policy, weighting is physics (§3.7,
  §8).
\item
  Erasure: derived-trace (\texttt{connect}) key problem stated as open
  (§6).
\item
  E3b (density) and E6 (reconcile) added; E3 baseline is timestamped
  RAG; E4 rewritten without S/P labels (§7).
\item
  Bench 2026-09-04 requalified as a loudness-readout pilot, not evidence
  for P1 (Availability).
\end{enumerate}

\end{document}
